\documentclass[a4paper]{article}

% \usepackage[default]{fontsetup}

\usepackage{fancyhdr}
\usepackage{extramarks}
\usepackage{amsmath}
\usepackage{amsthm}
\usepackage{amsfonts}
\usepackage{tikz}
\usepackage[plain]{algorithm}
\usepackage{algpseudocode}
\usepackage{enumerate}
\usepackage{tikz}

\usetikzlibrary{automata,positioning}

%
% Basic Document Settings
%  

\topmargin=-0.2in
\evensidemargin=0in
\oddsidemargin=0in
\textwidth=6.5in
\textheight=9.5in
\headsep=0.25in

\linespread{1.1}

\pagestyle{fancy}
\lhead{\hmwkAuthorName}
\chead{\hmwkClass : \hmwkTitle}
\rhead{\firstxmark}
\lfoot{\lastxmark}
\cfoot{\thepage}

\renewcommand\headrulewidth{0.4pt}
\renewcommand\footrulewidth{0.4pt}

\setlength\parindent{0pt}

%
% Create Problem Sections
%

\newcommand{\enterProblemHeader}[1]{
    \nobreak\extramarks{}{Problem \arabic{#1} continued on next page\ldots}\nobreak{}
    \nobreak\extramarks{Problem \arabic{#1} (continued)}{Problem \arabic{#1} continued on next page\ldots}\nobreak{}
}

\newcommand{\exitProblemHeader}[1]{
    \nobreak\extramarks{Problem \arabic{#1} (continued)}{Problem \arabic{#1} continued on next page\ldots}\nobreak{}
    \stepcounter{#1}
    \nobreak\extramarks{Problem \arabic{#1}}{}\nobreak{}
}

\newcommand*\circled[1]{\tikz[baseline=(char.base)]{
		\node[shape=circle,draw,inner sep=2pt] (char) {#1};}}


\setcounter{secnumdepth}{0}
\newcounter{partCounter}
\newcounter{homeworkProblemCounter}
\setcounter{homeworkProblemCounter}{1}
\nobreak\extramarks{Problem \arabic{homeworkProblemCounter}}{}\nobreak{}

%
% Homework Problem Environment
%
% This environment takes an optional argument. When given, it will adjust the
% problem counter. This is useful for when the problems given for your
% assignment aren't sequential. See the last 3 problems of this template for an
% example.
%

\newenvironment{homeworkProblem}[2][-1]{%
    \ifnum#1>0\relax
        \setcounter{homeworkProblemCounter}{#1}
    \fi
    \section{Problem \arabic{homeworkProblemCounter} (#2 points)}
    \setcounter{partCounter}{1}
    \enterProblemHeader{homeworkProblemCounter}
}{%
    \exitProblemHeader{homeworkProblemCounter}
}

%
% Homework Details
%   - Title
%   - Class
%   - Due date
%   - Name
%   - Student ID

\newcommand{\hmwkTitle}{Homework\ \#07}
\newcommand{\hmwkClass}{Probability \& Statistics for EECS}
\newcommand{\hmwkDueDate}{2026-11-8}
\newcommand{\hmwkAuthorName}{Jianxuan Yao}
\newcommand{\hmwkAuthorID}{2025521024}


%
% Title Page
%

\title{
    \vspace{2in}
    \textmd{\textbf{\hmwkClass:\\  \hmwkTitle}}\\
    \normalsize\vspace{0.1in}\small{Due\ on\ \hmwkDueDate\ at 23:59}\\
	\vspace{4in}
}

\author{
	Name: \textbf{\hmwkAuthorName} \\
	Student ID: \hmwkAuthorID}
\date{}

\renewcommand{\part}[1]{\textbf{\large Part \Alph{partCounter}}\stepcounter{partCounter}\\}

%
% Various Helper Commands
%

% Useful for algorithms
\newcommand{\alg}[1]{\textsc{\bfseries \footnotesize #1}}
% For derivatives
\newcommand{\deriv}[1]{\frac{\mathrm{d}}{\mathrm{d}x} (#1)}
% For partial derivatives
\newcommand{\pderiv}[2]{\frac{\partial}{\partial #1} (#2)}
% Integral dx
\newcommand{\dx}{\mathrm{d}x}
% Alias for the Solution section header
\newcommand{\solution}{\textbf{\large Solution}}
% Probability commands: Expectation, Variance, Covariance, Bias
\newcommand{\E}{\mathrm{E}}
\newcommand{\Var}{\mathrm{Var}}
\newcommand{\Cov}{\mathrm{Cov}}
\newcommand{\Bias}{\mathrm{Bias}}

\begin{document}


% \maketitle
% \thispagestyle{empty}
% \pagebreak

\date{
Due on Nov. 8, 2026, 23:59 UTC+8}
\title{SI 140A.02  Probability \& Statistics for EECS, Fall 2026 \\
Homework 7}
\maketitle
Read all the instructions below carefully before you start working on the assignment, and before you make a submission.
\begin{itemize}
    \item You are required to submit your homework in PDF format generated by \textbf{LaTeX}, otherwise you will get no points of this homework.
    \item You are required to write down all the major steps towards making your conclusions; otherwise you may obtain limited points of the problem.
    \item Write your homework in English; otherwise you will get no points of this homework.
    \item Any form of plagiarism will lead to $0$ point of this homework. 
\end{itemize}
\newpage
% Adapted from Joseph, Ch. 4, #24
\begin{homeworkProblem}{20}
Calvin and Hobbes play independent games. Calvin wins each game with probability $p$, and Hobbes wins with probability $q=1-p$, where $0<p<1$. The match ends as soon as either player has won two more games than the other.
\begin{enumerate}
\item[(a)] Let $T$ be the total number of games played. Find $\Pr(T=2k)$, $k=1,2,\ldots$. \emph{Hint:} Consider the games in consecutive pairs.
\item[(b)] Find $\E[T]$. Which value of $p$ maximizes this expectation?
\item[(c)] Find the probability that Calvin wins the match.
\end{enumerate}

\end{homeworkProblem}

\solution

\begin{enumerate}
    \item[(a)]
    
    Let Calvin win be $C$ and Hobbes win be $H$. The sequence of games can be represented as a string of $C$ and $H$, and $S$ be the sequence.

    $S$ must have $2k$ length, and the last two games must be either $CC$ or $HH$. The first $2k-2$ games is a joint of several $CH$ and $HC$.

    $$
    \Pr(CH) = \Pr(HC) = pq
    ,\quad
    \Pr(CC) = p^2
    ,\quad
    \Pr(HH) = q^2
    $$

    So $k \sim \mathrm{Geom}(2p^2 - 2p + 1)$.

    $$
    \Pr(T=2k) = (2p - 2p^2)^{k-1}(2p^2 - 2p + 1)
    $$

    \item[(b)]
    
    $$
    \E[k] = \frac{1}{2p^2 - 2p + 1}
    $$

    $$
    \E[T] = 2\E[k] = \frac{2}{2p^2 - 2p + 1}
    $$

    When $p = \frac{1}{2}$, $\E[T] = 4$ is the maximum.

    \item[(c)]
    
    Let $P_C$ be the probability of event "Calvin wins".

    $$
    P_C = p^2 + 2pq P_C
    $$

    $$
    P_C = \frac{p^2}{2p^2 - 2p + 1}
    $$
\end{enumerate}

\newpage
% Adapted from Joseph, Ch. 4, #40
\begin{homeworkProblem}{20}
A box contains $100$ shoelaces, each with two loose ends. Repeatedly pick two uniformly random loose ends and tie them together. A tie either joins two open pieces or closes one piece into a loop. Continue until all ends have been tied. Let $L$ be the total number of loops formed.
\begin{enumerate}
\item[(a)] When $m$ open pieces remain, find the probability that the next tie creates a loop.
\item[(b)] Find $\E[L]$ as a finite sum. Independence of the loop indicators is \emph{not} assumed.
\item[(c)] Find the probability that only one loop is formed by the end. A finite product is acceptable.
\end{enumerate}

\solution

\begin{enumerate}
    \item[(a)]
    
    $$
    \Pr(\text{next creats a loop} \mid m \text{ open pieces})
    = \frac{m}{\binom{2m}{2}} = \frac{1}{2m-1}
    $$

    \item[(b)]
    
    For every round, the shoelaces do either create a loop or not, and it can be seen as consistently diminishing one shoelace.

    The previous loop number does not affect the next round, so we can use the law of total expectation.

    Let $L_i$ be the loop number generated at $i$-th round. $L = L_1 + L_2 + \cdots + L_{100}$.

    $$
    \E[L_i] = \frac{1}{2(100-i+1)-1}
    $$

    $$
    \E[L]
    = \sum_{i=1}^{100} \E[L_i]
    = \sum_{i=1}^{100} \frac{1}{2(100-i+1)-1}
    = \sum_{i=1}^{100} \frac{1}{201-2i}
    = \sum_{i=1}^{100} \frac{1}{2i-1}
    $$

    \item[(c)]
    
    When there only remains one shoelace, the last tie must create a loop. So the previous $99$ ties must not create a loop.

    $$
    \Pr(L=1)
    = \prod_{i=1}^{99} \left(1 - \frac{1}{2(100-i+1)-1}\right)
    = \prod_{i=1}^{99} \frac{2(100-i+1)-2}{2(100-i+1)-1}
    = \prod_{i=1}^{99} \frac{2i-1}{2i}
    $$
\end{enumerate}

\end{homeworkProblem}
\newpage
% Adapted from Joseph, Ch. 4, #28 (zero-inflated Poisson)
\begin{homeworkProblem}{20}
A monitoring device is \emph{inactive} on a given day with probability $\alpha\in(0,1)$. On an active day, each of $n$ independent sensors produces an alert with a small probability $u$. Approximate the active-day alert count by $\mathrm{Pois}(\lambda)$, where $\lambda=nu$ (with large $n$ and small $u$). Let $X$ be the number of alerts on a randomly selected day. The inactive/active status is independent of the sensor outcomes.
\begin{enumerate}
\item[(a)] Identify the original distribution of the active-day alert count before approximation. Under the Poisson approximation, find $\Pr(X=0)$ and $\Pr(X=k)$ for all integers $k\ge1$.
\item[(b)] Find $\E[X]$ and $\Var(X)$. You may introduce independent variables $I\sim\mathrm{Bern}(\alpha)$ and $Y\sim\mathrm{Pois}(\lambda)$ with $X=(1-I)Y$.
\item[(c)] Take $n=200$, $u=0.01$ and $\alpha=1/5$ (so $\lambda=2$). Compare $\Pr(X=0)$ under the mixture model with the zero probability under $\mathrm{Pois}(\E[X])$. Explain briefly why using only the mean to select a Poisson model can be misleading.
\end{enumerate}

\end{homeworkProblem}

\solution

\begin{enumerate}
    \item[(a)]
    
    The orginal distribution of the active-day alert count is $\mathrm{Bin}(n,u)$.

    Let $A$ be the event that the day is active.

    $\Pr(A) = 1 - \alpha$.

    $$
    \begin{aligned}
        \Pr(X = 0)
        &= \Pr(X = 0 \mid A^c)\Pr(A^c) + \Pr(X = 0 \mid A)\Pr(A) \\
        &= 1 \cdot \alpha + e^{-\lambda} \cdot (1 - \alpha) \\
        &= \alpha + (1 - \alpha)e^{-\lambda}
    \end{aligned}
    $$

    When $k \ge 1$,
    
    $$
    \begin{aligned}
        \Pr(X = k)
        &= \Pr(X = k \mid A^c)\Pr(A^c) + \Pr(X = k \mid A)\Pr(A) \\
        &= 0 \cdot \alpha + \frac{\lambda^k e^{-\lambda}}{k!} \cdot (1 - \alpha) \\
        &= (1 - \alpha)\frac{\lambda^k e^{-\lambda}}{k!}
    \end{aligned}
    $$

    \item[(b)]
    
    Let $I$ be the indicator of the event that the day is inactive, and $Y$ be the number of alerts on an active day. Then $I \sim \mathrm{Bern}(\alpha)$ and $Y \sim \mathrm{Pois}(\lambda)$, and $X = (1 - I)Y$.

    $I$ and $Y$ are independent.

    $$
    \E[1 - I] = 1 - \alpha, \quad \Var(1 - I) = \alpha(1 - \alpha)
    $$

    $$
    \E[Y] = \lambda, \quad \Var(Y) = \lambda, \quad \E[Y^2] = \Var(Y) + (\E[Y])^2 = \lambda + \lambda^2
    $$

    $$
    \E[X] = \E[(1 - I)Y] = \E[1 - I] \cdot \E[Y] = (1 - \alpha)\lambda
    $$

    $$
    \E[X^2] = \E[(1 - I)^2 Y^2] = \E[(1 - I) Y^2] = \E[1 - I] \cdot \E[Y^2] = (1 - \alpha)(\lambda + \lambda^2)
    $$

    $$
    \begin{aligned}
        \Var(X)
        &= \E[X^2] - (\E[X])^2 \\ &= (1 - \alpha)(\lambda + \lambda^2) - (1 - \alpha)^2 \lambda^2 \\
        &= (1 - \alpha)\lambda(1 + \alpha\lambda)
    \end{aligned}
    $$

    \item[(c)]
    
    $$
    \E[X] = (1 - 0.2) \times 2 = 0.8 \times 2 = 1.6
    $$

    $$
    \Pr(X = 0) = \alpha + (1 - \alpha)e^{-\lambda} = 0.2 + 0.8 e^{-2} \approx 30.83\%
    $$

    Let $W \sim \text{Pois}(1.6)$,

    $$
    \Pr(W = 0) = e^{-1.6} \approx 20.19\%
    $$

    In this case, the mixture distribution has an additional variance due to the switch between active and inactive. As a result, its variance is larger than expectation, leading to misleading results.
\end{enumerate}
\newpage
% Original: overlapping patterns and second moments
\begin{homeworkProblem}{20}
A fair coin is tossed independently $n\ge5$ times. Let $X$ count the occurrences of the pattern HTH in consecutive positions, allowing overlaps (e.g., HTHTH contains two occurrences). For $i=1,\ldots,n-2$, let $I_i$ indicate the occurrence starting at position $i$, so $X=\sum_{i=1}^{n-2}I_i$.
\begin{enumerate}
\item[(a)] Find $\E[X]$.
\item[(b)] Compute $\Pr(I_i=I_{i+1}=1)$ and $\Pr(I_i=I_{i+2}=1)$. Explain why $I_i$ and $I_j$ are independent whenever $|i-j|\ge3$.
\item[(c)] Derive $\Var(X)$ as a function of $n$. Evaluate both $\E[X]$ and $\Var(X)$ when $n=20$.
\end{enumerate}

\end{homeworkProblem}

\solution
\begin{enumerate}
    \item[(a)]
    
    $$
    \E[X] = \sum_{i=1}^{n-2} \E[I_i]
    $$

    $$
    \E[I_i] = \frac{1}{8}
    $$

    $$
    \E[X] = \frac{n-2}{8}
    $$

    \item[(b)]
    
    $$
    \Pr(I_i=I_{i+1}=1) = 0
    $$

    $$
    \Pr(I_i=I_{i+2}=1) = \frac{1}{32}
    $$

    $I_i$ and $I_j$ are independent whenever $|i-j|\ge3$ because they do not overlap, and the coin tosses are independent.

    \item[(c)]
    
    $$
    \Pr(I_i I_j)
    =
    \begin{cases}
        0, & \vert i-j \vert = 1 \\
        \frac{1}{32}, & \vert i-j \vert = 2 \\
        \frac{1}{64}, & \vert i-j \vert \ge 3
    \end{cases}
    $$

    $$
    X^2 = \sum_{i=1}^{n-2} I_i^2 + 2\sum_{1\le i<j\le n-2} I_i I_j
    $$

    $$
    \begin{aligned}
    \E[X^2]
    &= \sum_{i=1}^{n-2} \E[I_i^2] + 2\sum_{1\le i<j\le n-2} \E[I_i I_j] \\
    &= \sum_{i=1}^{n-2} \E[I_i] + 2\sum_{1\le i<j\le n-2} \E[I_i I_j] \\
    &=
    \sum_{i=1}^{n-2} \E[I_i]
    + 2\sum_{i=1}^{n-3} \E[I_i I_{i+1}]
    + 2\sum_{i=1}^{n-4} \E[I_i I_{i+2}]
    + 2\sum_{i=1}^{n-5} \sum_{j=i+3}^{n-2} \E[I_i I_j] \\
    &= \sum_{i=1}^{n-2} \frac{1}{8}
    + 2\sum_{i=1}^{n-3} 0
    + 2\sum_{i=1}^{n-4} \frac{1}{32}
    + 2\sum_{i=1}^{n-5} \sum_{j=i+3}^{n-2} \frac{1}{64} \\
    &= \frac{n-2}{8} + \frac{n-4}{16} + \frac{(n-5)(n-4)}{64} \\
    &= \frac{n^2+3n-12}{64}
    \end{aligned}
    $$

    $$
    \Var(X) = \E[X^2] - (\E[X])^2 = \frac{7n - 16}{64}
    $$

    When $n = 20$, $\E[X] = \frac{9}{4}$ and $\Var(X) = \frac{31}{16}$.
\end{enumerate}

\newpage
% Original: negative binomial, binomial, PGF and moments
\begin{homeworkProblem}{20}
A transmitter repeats independent attempts, each successful with probability $p\in(0,1)$. It stops at the \emph{third} successful attempt. Let $N$ be its total number of attempts, and $Y=N-3$ the number of failures before the third success.
\begin{enumerate}
\item[(a)] Find the PMF of $Y$ and identify its distribution.
\item[(b)] Express $\Pr(N\le6)$ in terms of a binomial random variable counting successes in the first six attempts. Evaluate this probability when $p=1/2$.
\item[(c)] View $Y$ as a sum of three independent geometric random variables, each counting failures before a success. Find its PGF $G_Y(t)=\E[t^Y]$ and use $G_Y'(1)$ and $G_Y''(1)$ to obtain $\E[N]$ and $\Var(N)$.
\end{enumerate}

\end{homeworkProblem}

\solution
\begin{enumerate}
    \item[(a)]

    $$
    \Pr(N = k) = \binom{k-1}{2} p^3 (1-p)^{k-3}, \quad k = 3,4,\ldots
    $$

    $$
    \Pr(Y = k) = \binom{k+2}{2} p^3 (1-p)^k, \quad k = 0,1,\ldots
    $$

    $$
    Y \sim \mathrm{NBin}(3,p)
    $$

    \item[(b)]
    
    Let $X$ be the number of successes in the first six attempts, so $X \sim \text{Binomial}(6, p)$.

    $$
    \Pr(N \le 6) = \Pr(X \ge 3) = \sum_{k=3}^6 \binom{6}{k} p^k (1 - p)^{6 - k}
    $$

    When $p = 1/2$,

    $$
    \Pr(N \le 6) = \frac{21}{32}
    $$

    \item[(c)]
    
    $$
    Y = X_1 + X_2 + X_3
    $$

    $$
    X_i \sim \text{Geom}(p), \quad i = 1,2,3
    $$

    $$
    \Pr(X_i = k) = (1 - p)^k p, \quad k = 0, 1, 2, \dots
    $$

    $$
    G_X(t)
    = \E[t^{X_i}]
    = \sum_{k=0}^\infty t^k (1 - p)^k p
    = p \sum_{k=0}^\infty [(1 - p)t]^k
    = \frac{p}{1 - (1 - p)t}, \quad \text{where } \vert{}t\vert{} < \frac{1}{1 - p}
    $$

    $X_1, X_2, X_3$ are independent, so

    $$
    G_Y(t)
    = \E[t^Y]
    = [G_X(t)]^3
    = \left( \frac{p}{1 - (1 - p)t} \right)^3
    $$

    $$
    G_Y'(t) = 3 p^3 q (1 - qt)^{-4}
    $$

    $$
    G_Y''(t) = 12 p^3 q^2 (1 - qt)^{-5}
    $$

    $$
    G_Y'(1) = 3 p^3 q (1 - q)^{-4} = \frac{3(1-p)}{p}
    $$

    $$
    G_Y''(1) = 12 p^3 q^2 (1 - q)^{-5} = \frac{12(1-p)^2}{p^2}
    $$

    $$
    \E[Y] = G_Y'(1) = \frac{3(1-p)}{p}
    $$

    $$
    \E[N] = \E[Y] + 3 = \frac{3}{p}
    $$

    $$
    \E[Y^2] - \E[Y] = G_Y''(1) = \frac{12(1-p)^2}{p^2}
    $$

    $$
    \E[(N-3)^2] - \E[N-3] = \frac{12(1-p)^2}{p^2}
    $$

    $$
    \E[N^2] - 7 \E[N] + 12 = \frac{12(1-p)^2}{p^2}
    $$

    $$
    \E[N^2] = \frac{12(1-p)^2}{p^2} + 7 \E[N] - 12 = \frac{12(1-p)^2}{p^2} + \frac{21}{p} - 12
    $$

    $$
    \Var(N) = \E[N^2] - (\E[N])^2 = \frac{12(1-p)^2}{p^2} + \frac{21}{p} - 12 - \frac{9}{p^2} = \frac{3(1-p)}{p^2}
    $$
\end{enumerate}
\end{document}